% GENERATED FILE. Edit canonical lessons, then run npm run generate:books.
% canonical-source-hash: fnv1a64:1ebda8ea664f8735
% canonical-lessons: JA-C342-gen-in, JA-C342-eikyou, JA-C342-kekka, JA-C342-nazenara, JA-C342-taido

\chapter{Causes and Effects}
\label{ch:ja-causes-and-effects}
\hlchaptermodality{342}

\begin{chapteropening}
\textbf{By the end of this chapter:} \emph{I can say what caused something and what came of it.}

Complete the last lesson of chapter 342: I can say what caused something and what came of it.
\end{chapteropening}

\section[\texorpdfstring{\ja{げんいん} (gen'in)}{げんいん (gen'in)}]{\ja{げんいん} (gen'in) — a cause}

\label{lesson:JA-C342-gen-in}



\begin{quote}
Before the new one: say the Japanese for as much as possible, then the Japanese for for example.
\end{quote}

\subsection*{You'll want to know: \ja{げんいん}}
\textbf{\ja{げんいん}} — \emph{gen'in} — \textquotedblleft{}a cause\textquotedblright{}.

\begin{quote}
\textbf{\ja{じこの} \ja{げんいんが} \ja{わからないので、しらべています。}} — \emph{jiko no gen'in ga wakaranai node, shirabete imasu} — \textquotedblleft{}Nobody knows the cause of the accident, so it is being investigated.\textquotedblright{}
\end{quote}

\begin{cousinweb}
\textbf{\ja{げんいん}} (\emph{gen'in}) is Sino-Japanese (kango): two Chinese characters read with their on'yomi, readings borrowed from Chinese. They are \emph{gen}, \textquotedblleft{}original\textquotedblright{} and \emph{in}, \textquotedblleft{}cause\textquotedblright{} (Middle Chinese \emph{'jin}) (KANJIDIC2; Middle Chinese in Baxter's transcription).
\end{cousinweb}

\begin{grammarlens}[title={What you've built}]
One more word for giving reasons and opinions.
\end{grammarlens}

\subsection*{Your turn}
\begin{itemize}
\raggedright
  \item \emph{Say it:} \emph{gen'in}
  \item \emph{Say it:} \emph{gen'in}, once more
  \item \emph{Say it:} \emph{jiko no gen'in ga wakaranai node, shirabete imasu} — Nobody knows the cause of the accident, so it is being investigated.
  \item \emph{Recall:} say the Japanese for as much as possible, then the Japanese for for example, then say \emph{gen'in} again
\end{itemize}

\begin{tcolorbox}[breakable,colback=teal!4,colframe=teal!35!black,title={Before you move on}]
What does \ja{げんいん} mean? (\textquotedblleft{}A cause\textquotedblright{}.) Say it once more.
\end{tcolorbox}
\section[\texorpdfstring{\ja{えいきょう} (eikyō)}{えいきょう (eikyō)}]{\ja{えいきょう} (eikyō) — an influence; an effect}

\label{lesson:JA-C342-eikyou}



\begin{quote}
Before the new one: say the Japanese for for example, then the Japanese for a cause.
\end{quote}

\subsection*{You'll want to know: \ja{えいきょう}}
\textbf{\ja{えいきょう}} — \emph{eikyō} — \textquotedblleft{}an influence; an effect\textquotedblright{}.

\begin{quote}
\textbf{\ja{たいふうの} \ja{えいきょうで、でんしゃが} \ja{とまりました。}} — \emph{taifū no eikyō de, densha ga tomarimashita} — \textquotedblleft{}The trains stopped because of the typhoon.\textquotedblright{}
\end{quote}

\begin{cousinweb}
\textbf{\ja{えいきょう}} (\emph{eikyō}) is Sino-Japanese (kango): two Chinese characters read with their on'yomi, readings borrowed from Chinese. They are \emph{ei}, \textquotedblleft{}shadow\textquotedblright{} (Middle Chinese \emph{'jaengX}) and \emph{kyō}, \textquotedblleft{}echo\textquotedblright{} (\emph{xjangX}) (KANJIDIC2; Middle Chinese in Baxter's transcription).
\end{cousinweb}

\begin{grammarlens}[title={What you've built}]
One more word for giving reasons and opinions.
\end{grammarlens}

\subsection*{Your turn}
\begin{itemize}
\raggedright
  \item \emph{Say it:} \emph{eikyō}
  \item \emph{Say it:} \emph{eikyō}, once more
  \item \emph{Say it:} \emph{taifū no eikyō de, densha ga tomarimashita} — The trains stopped because of the typhoon.
  \item \emph{Recall:} say the Japanese for for example, then the Japanese for a cause, then say \emph{eikyō} again
\end{itemize}

\begin{tcolorbox}[breakable,colback=teal!4,colframe=teal!35!black,title={Before you move on}]
What does \ja{えいきょう} mean? (\textquotedblleft{}An influence; an effect\textquotedblright{}.) Say it once more.
\end{tcolorbox}
\section[\texorpdfstring{\ja{けっか} (kekka)}{けっか (kekka)}]{\ja{けっか} (kekka) — a result; an outcome}

\label{lesson:JA-C342-kekka}



\begin{quote}
Before the new one: say the Japanese for a cause, then the Japanese for an influence.
\end{quote}

\subsection*{You'll want to know: \ja{けっか}}
\textbf{\ja{けっか}} — \emph{kekka} — \textquotedblleft{}a result; an outcome\textquotedblright{}.

\begin{quote}
\textbf{\ja{がんばったので、いい} \ja{けっかに} \ja{なりました。}} — \emph{ganbatta node, ii kekka ni narimashita} — \textquotedblleft{}We worked hard, so it turned out well.\textquotedblright{}
\end{quote}

\begin{cousinweb}
\textbf{\ja{けっか}} (\emph{kekka}) is Sino-Japanese (kango): two Chinese characters read with their on'yomi, readings borrowed from Chinese. They are \emph{ketsu}, \textquotedblleft{}tie\textquotedblright{} (Middle Chinese \emph{ket}) and \emph{ka}, \textquotedblleft{}fruit\textquotedblright{} (KANJIDIC2; Middle Chinese in Baxter's transcription). The two join as \emph{kek-ka}.
\end{cousinweb}

\begin{grammarlens}[title={What you've built}]
One more word for giving reasons and opinions.
\end{grammarlens}

\subsection*{Your turn}
\begin{itemize}
\raggedright
  \item \emph{Say it:} \emph{kekka}
  \item \emph{Say it:} \emph{kekka}, once more
  \item \emph{Say it:} \emph{ganbatta node, ii kekka ni narimashita} — We worked hard, so it turned out well.
  \item \emph{Recall:} say the Japanese for a cause, then the Japanese for an influence, then say \emph{kekka} again
\end{itemize}

\begin{tcolorbox}[breakable,colback=teal!4,colframe=teal!35!black,title={Before you move on}]
What does \ja{けっか} mean? (\textquotedblleft{}A result; an outcome\textquotedblright{}.) Say it once more.
\end{tcolorbox}
\section[\texorpdfstring{\ja{なぜなら} (nazenara)}{なぜなら (nazenara)}]{\ja{なぜなら} (nazenara) — because (the reason is)}

\label{lesson:JA-C342-nazenara}



\begin{quote}
Before the new one: say the Japanese for an influence, then the Japanese for a result.
\end{quote}

\subsection*{You'll want to know: \ja{なぜなら}}
\textbf{\ja{なぜなら}} — \emph{nazenara} — \textquotedblleft{}because (the reason is)\textquotedblright{}.

\begin{quote}
\textbf{\ja{なぜなら、あした} \ja{しけんが} \ja{あるからです。}} — \emph{nazenara, ashita shiken ga aru kara desu} — \textquotedblleft{}The reason is that I have an exam tomorrow.\textquotedblright{}
\end{quote}

\begin{cousinweb}
\textbf{\ja{なぜなら}} (\emph{nazenara}) is a native Japanese word (wago), inherited rather than borrowed from Chinese: it is built on \emph{naze}, \textquotedblleft{}why\textquotedblright{}, and JMdict lists it as a conjunction, usually written in kana (JMdict).
\end{cousinweb}

\begin{grammarlens}[title={What you've built}]
One more word for giving reasons and opinions.
\end{grammarlens}

\subsection*{Your turn}
\begin{itemize}
\raggedright
  \item \emph{Say it:} \emph{nazenara}
  \item \emph{Say it:} \emph{nazenara}, once more
  \item \emph{Say it:} \emph{nazenara, ashita shiken ga aru kara desu} — The reason is that I have an exam tomorrow.
  \item \emph{Recall:} say the Japanese for an influence, then the Japanese for a result, then say \emph{nazenara} again
\end{itemize}

\begin{tcolorbox}[breakable,colback=teal!4,colframe=teal!35!black,title={Before you move on}]
What does \ja{なぜなら} mean? (\textquotedblleft{}Because (the reason is)\textquotedblright{}.) Say it once more.
\end{tcolorbox}
\section[\texorpdfstring{\ja{たいど} (taido)}{たいど (taido)}]{\ja{たいど} (taido) — an attitude; a manner}

\label{lesson:JA-C342-taido}



\begin{quote}
Before the new one: say the Japanese for a result, then the Japanese for the reason is.
\end{quote}

\subsection*{You'll want to know: \ja{たいど}}
\textbf{\ja{たいど}} — \emph{taido} — \textquotedblleft{}an attitude; a manner\textquotedblright{}.

\begin{quote}
\textbf{\ja{てんいんの} \ja{たいどが} \ja{きらいだから、あの} \ja{みせに} \ja{いきません。}} — \emph{ten'in no taido ga kirai da kara, ano mise ni ikimasen} — \textquotedblleft{}I dislike the shop assistant's attitude, so I don't go to that shop.\textquotedblright{}
\end{quote}

\begin{cousinweb}
\textbf{\ja{たいど}} (\emph{taido}) is Sino-Japanese (kango): two Chinese characters read with their on'yomi, readings borrowed from Chinese. They are \emph{tai}, \textquotedblleft{}attitude\textquotedblright{} (Middle Chinese \emph{thojH}) and \emph{do}, \textquotedblleft{}degrees\textquotedblright{} (\emph{duH}) (KANJIDIC2; Middle Chinese in Baxter's transcription).
\end{cousinweb}

\begin{grammarlens}[title={What you've built}]
One more word for giving reasons and opinions.
\end{grammarlens}

\subsection*{Your turn}
\begin{itemize}
\raggedright
  \item \emph{Say it:} \emph{taido}
  \item \emph{Say it:} \emph{taido}, once more
  \item \emph{Say it:} \emph{ten'in no taido ga kirai da kara, ano mise ni ikimasen} — I dislike the shop assistant's attitude, so I don't go to that shop.
  \item \emph{Recall:} say the Japanese for a result, then the Japanese for the reason is, then say \emph{taido} again
\end{itemize}

\begin{tcolorbox}[breakable,colback=teal!4,colframe=teal!35!black,title={Before you move on}]
What does \ja{たいど} mean? (\textquotedblleft{}An attitude; a manner\textquotedblright{}.) Say it once more.
\end{tcolorbox}
